\documentclass[../../../include/open-logic-section]{subfiles}

\begin{document}

\olfileid{sfr}{infinite}{card-sb}

\olsection{Bijlage: de stelling van Schr\"oder-Bernstein bewijzen}

Voordat we de na\"ieve verzamelingenleer achter ons laten, bekijken we
nog één laatste na\"ief bewijs. Het is wel behoorlijk slim opgebouwd!
Het gaat om een bewijs van de stelling van Schr\"oder-Bernstein
(\olref[sfr][siz][sb]{thm:schroder-bernstein}): als
$\cardle{A}{B}$ en $\cardle{B}{A}$, dan $\cardeq{A}{B}$. Anders
gezegd: bij injecties $f \colon A \to B$ en $g \colon B \to A$
bestaat er een bijectie $h \colon A \to B$.

In dit hoofdstuk gebruikten we Dedekinds idee van \emph{afsluitingen}.
Met dat idee gaf Dedekind een mooi bewijs van de stelling van
Schr\"oder-Bernstein. Dat bewijs staat hier. Onze versie volgt
\citet[pp.~157--8]{Potter2004} bijna helemaal. Daar staat een iets
andere aanpak die in de kern hetzelfde is. Als je even googelt, zie je
ook dat voor deze stelling---net als voor de irrationaliteit van
$\sqrt{2}$---\emph{veel} interessante en heel verschillende bewijzen
bestaan.

We gebruiken bijna dezelfde notatie als in
\olref[sfr][infinite][dedekind]{Closure} en schrijven
\[
\Closureofunder{f}{B} = \bigcap \Setabs{X}{B \subseteq X
\text{ en $X$ is $f$-gesloten}}
\]
voor elke verzameling~$B$ en functie~$f$. De verzameling
$\Closureofunder{f}{B}$ is de kleinste $f$-gesloten verzameling die~$B$
bevat. Dat betekent het volgende:

\begin{lem}\ollabel{Closureprops}
	Voor elke functie $f$ en elke $B$:
	\begin{enumerate}
		\item\ollabel{Closurehaselem} $B \subseteq \Closureofunder{f}{B}$; en
		\item\ollabel{Closureclosed} $\Closureofunder{f}{B}$ is $f$-gesloten; en
		\item\ollabel{Closuresmallest} als $X$ $f$-gesloten is en $B
		\subseteq X$, dan $\Closureofunder{f}{B} \subseteq X$.
	\end{enumerate}
\end{lem}

\begin{proof}
Dit gaat precies zoals in
\olref[sfr][infinite][dedekind]{closureproperties}.
\end{proof}

Voor de stelling van Schr\"oder-Bernstein hebben we nog één laatste feit
nodig:

\begin{prop}\ollabel{sbhelper}
Als $A \subseteq B \subseteq C$ en $A \approx C$, dan $\cardeq{\cardeq{A}{B}}{C}$.
\end{prop}

\begin{proof}
Neem een bijectie $f \colon C \to A$. Stel $F =
\Closureofunder{f}{C \setminus B}$ en definieer als volgt een functie
$g$ met domein~$C$:
\[
	g(x) =
	\begin{cases}
			f(x) &\text{als $x \in F$}\\
			x & \text{anders}
		\end{cases}
\]
We laten zien dat $g$ een bijectie van $C \to B$ is. Dan is ook
$\comp{f^{-1}}{g} \colon A \to B$ een bijectie, en is het bewijs
klaar.

Eerst beweren we: als $x \in F$ maar $y\notin F$, dan $g(x) \neq g(y)$.
Stel, om een tegenspraak te krijgen, dat ze toch gelijk zijn. Dan geldt
$y = g(y) = g(x) = f(x)$. Omdat $x \in F$ en $F$ volgens
\olref{Closureprops} $f$-gesloten is, geldt $y = f(x) \in F$. Dat is
een tegenspraak.

Stel nu dat $g(x) = g(y)$. Uit wat we net bewezen volgt: $x \in F$ dan
en slechts dan als $y \in F$. Als $x, y \in F$, dan is $f(x) = g (x)
= g(y) = f(y)$. Dus $x = y$, want $f$ is een bijectie. Als
$x, y \notin F$, dan $x = g(x) = g(y) = y$. De functie~$g$ is dus
injectief.

We moeten nog bewijzen dat $\ran{g} = B$. Eerst laten we zien dat
$\ran{g} \subseteq B$. Neem $z \in C$. Als $z \in F$, dan is
$g(z) = f(z) \in A \subseteq B$. Als $z \notin F$, dan geldt
$z \notin C \setminus B$, want $C \setminus B \subseteq F$. Dus geldt
$z \in B$ en $g(z) = z \in B$. Daarmee is
$\ran{g} \subseteq B$. Nu de andere kant. Neem $x \in B \subseteq C$.
Als $x \notin F$, dan $g(x) = x$.
Als $x \in F$, dan is $x = f(y)$ voor een $y \in F$. Anders zou
$F \setminus \{x\}$ namelijk $f$-gesloten zijn en $C\setminus B$
bevatten. Volgens \olref{Closureprops} kan dat niet. Nu geldt
$g(y) = f(y) = x$.
\end{proof}

Nu volgt het bewijs van de hoofdstelling. Voor een functie~$h$ en een
verzameling~$D$ hadden we $\funimage{h}{D} = \Setabs{h(x)}{x \in D}$
gedefinieerd.

\begin{proof}[Bewijs van de stelling van Schr\"oder-Bernstein] Neem
injecties $f \colon A \to B$ en $g \colon B \to A$. Omdat
$\funimage{f}{A} \subseteq B$, geldt
$\funimage{g}{\funimage{f}{A}} \subseteq g[B] \subseteq A$. Ook is
$\comp{f}{g} \colon A \to \funimage{g}{\funimage{f}{A}}$ injectief,
want $g$ en $f$ zijn dat allebei. Deze functie is zelfs bijectief: we
hebben het codomein van $\comp{f}{g}$ precies gelijk aan het beeld van
die functie gekozen. Dus
$\cardeq{\funimage{g}{\funimage{f}{A}}}{A}$. Uit
\olref{sbhelper} volgt daarom dat er een bijectie $h \colon A \to
\funimage{g}{B}$ bestaat. Ook is $g^{-1}$ een bijectie van
$\funimage{g}{B} \to B$. Dus is $\comp{h}{g^{-1}} \colon A \to B$
een bijectie.
\end{proof}
\end{document}